% =============================================================
% Chapter 2 --- LLM Fundamentals
% =============================================================

\chapter{LLM Fundamentals}
\label{ch:fundamentals}

\todo{This chapter is a stub. It will introduce the building blocks
that the rest of the book takes for granted: the autoregressive
generation loop, tokens versus characters, the chat-message
abstraction (system / user / assistant turns), temperature and
sampling, and \emph{structured outputs}. The structured-outputs
section is the prerequisite for understanding tool calls
(Chapter~\ref{ch:toolcalls}), and is sketched briefly below so the
forward references compile.}

\section{Structured Outputs}
\label{sec:structured}

By default a language model emits a stream of tokens that decode to
free-form natural language. Many downstream uses need more: a JSON
object with a fixed shape, an XML document, a SQL query, a function
signature. Three families of techniques are in common use, in
increasing order of guarantees:

\begin{enumerate}[label=(\arabic*)]
\item \textbf{Prompting.} The system prompt instructs the model to
  emit, say, ``a single JSON object on one line, no prose, no code
  fences''. This works most of the time on a sufficiently large
  model and almost always with one or two few-shot examples. It is
  the cheapest option and the one most readers will start with.
\item \textbf{Constrained decoding.} The runtime intercepts the
  token stream and masks out any token that would violate a target
  grammar. Implementations include \code{llama.cpp}'s GBNF grammars,
  \code{Outlines}, \code{lm-format-enforcer}, and the structured-output
  modes shipped by hosted providers. The output is \emph{guaranteed}
  to parse against the schema.
\item \textbf{Fine-tuning.} The model is trained on examples of the
  desired structured format. Used in concert with constrained
  decoding it produces the most reliable output, at the cost of a
  training run.
\end{enumerate}

\noindent
A tool call is the canonical application of structured outputs: a
JSON object with a known schema, intercepted by the runtime,
dispatched, and answered. Chapter~\ref{ch:toolcalls} treats the tool
call as the unit of work and assumes structured-output reliability
is a solved problem at the layer below.

\todo{Expand: a worked example of constrained decoding against a
local model; a comparison of the schema languages used by major
providers (OpenAI's tool schema, Anthropic's, Gemma 4's auto-derived
schema from Python type hints); the trade-off between
prompt-engineering reliability and constrained-decoding rigidity.}
